\titlespacing*{\chapter}{0pt}{-2cm}{30pt}
\chapter*{GENERAL INTRODUCTION}
\titlespacing*{\chapter}{0pt}{0pt}{50pt}
\markboth{\MakeUppercase{GENERAL INTRODUCTION}}{}
\addcontentsline{toc}{chapter}{GENERAL INTRODUCTION}
\adjustmtc
\thispagestyle{MyStyle}

Organizational practices are now dependent not only on human resources but also on the ability to manage processes effectively and adapt to rapid change driven by the evolution of digital technologies. In this context, recruiting is strategically important, as it directly impacts team building and organizational success. However, the complexity of this process has increased because of the growing number of applications and the demand for faster, more accurate treatment. Consequently, organizations are turning to digital solutions that centralize information, improve workflows, and support decision-making, turning recruitment into a structured process that requires efficient tools.

\par

In this setting, Sopra HR Software is recognized as a leading player in human resource management solutions. Operating in a field where technological innovation and service quality are essential, the company emphasizes the need to improve its internal processes and business tools. In this sphere, the digitalization of recruitment is a major lever for streamlining application processing, enhancing the user experience, and supporting recruiters in their day-to-day activities.

\par

The implementation of an intelligent recruitment portal aims to meet both operational and strategic needs by providing a more cohesive framework for managing job offers, applications, and interactions among all stakeholders involved in recruitment. Consequently, a central question arises: how can a recruitment portal be designed to centralize applications, facilitate CV processing, and effectively support recruiters in their decision-making, while aligning with a logic of continuous improvement in the recruitment process?

\par

To explore the identified problem and present the work conducted during the internship, the report is divided into three main chapters:

\begin{itemize}[label={-}]
  \item \textbf{Chapter One} presents the general framework of the project by introducing the host organization and situating the project within its context. It analyzes the existing recruitment process and its limitations, compares recruitment solutions, and justifies the proposed intelligent recruitment portal.

  \item \textbf{Chapter Two} addresses requirements analysis and system design. It describes the functional and non-functional requirements, highlights the expected features, and introduces the main design components used to build the portal.

  \item \textbf{Chapter Three} is devoted to the development and implementation phase. It outlines the steps involved in building the application, details the implemented features, and summarizes the main outcomes achieved.
\end{itemize}

\par

Finally, the report concludes with a summary that recaps the context of the work, reviews the project, and presents possible directions for future enhancements.
